\subsection{SAT formulation and implementation}
\label{app:sat-primer}

Boolean satisfiability (SAT) asks whether binary variables can be
assigned values that make a logical formula true. The values 1 and 0
represent true and false. The operations $\wedge$, $\vee$, and $\neg$
denote AND, OR, and NOT, respectively. Thus, $x\wedge y$ is true when
both variables are true, $x\vee y$ is true when at least one is true,
and $\neg x$ is true when $x$ is false. A variable or its negation is
called a \emph{literal}. A \emph{clause} is a disjunction of literals,
and a formula in \emph{conjunctive normal form} (CNF) is a conjunction
of clauses. Such a formula is true precisely when every clause is true.
For example, $(x_1\vee x_2)\wedge\neg x_1$ is satisfied by $x_1=0$
and $x_2=1$. A formula is \emph{satisfiable} if at least one satisfying
assignment exists and \emph{unsatisfiable} otherwise. See
\citet{biere2021handbook} for further background.

In our problem, $\QQ$ is fixed, and the unknowns are the entries of
$\bar{\QQ}$ and $\HH$. Under the assumptions of
Proposition~\ref{prop:sat-nonequivalence}, the algebraic condition
fails precisely when these matrices satisfy~\eqref{eq:sat_cond}.
The binary variables $x_{j\ell}$ and $h_{\ell k}$ represent
$\bar q_{j\ell}$ and the entries of $\HH$, respectively. We give an
explicit CNF encoding of~\eqref{eq:sat-factorization} below. Here $J$
denotes the number of rows of the matrix supplied to the encoding;
for Algorithm~\ref{alg:id}, this is the number of rows after basis
reduction.

First consider the Boolean factorization. If $q_{jk}=0$, every term
$x_{j\ell}\wedge h_{\ell k}$ must be zero. This is enforced by the
$K$ clauses
\begin{equation}\label{eq:supp-sat-zero}
\neg x_{j\ell}\vee\neg h_{\ell k},
\qquad \ell\in[K].
\end{equation}
If $q_{jk}=1$, at least one of these terms must be one. For this entry,
introduce auxiliary binary variables $z_{j\ell k}$, $\ell\in[K]$,
that select a term whose two factors are one. We impose
$z_{j\ell k}\Rightarrow(x_{j\ell}\wedge h_{\ell k})$ using
\begin{equation}\label{eq:supp-sat-and}
\neg z_{j\ell k}\vee x_{j\ell},
\qquad
\neg z_{j\ell k}\vee h_{\ell k},
\end{equation}
and add the clause
\begin{equation}\label{eq:supp-sat-positive}
\bigvee_{\ell=1}^K z_{j\ell k}.
\end{equation}
The reverse implication is unnecessary: the $z$-variables need only
witness one valid term. For any assignment satisfying the Boolean
factorization, such a term can be selected for every positive entry.
Conversely, a satisfying assignment of these clauses supplies such a
term and forbids every positive term at a zero entry. Thus, applying
these constraints to every entry of $\QQ$ gives exactly the Boolean
factorization in~\eqref{eq:sat-factorization}, after existentially
quantifying the auxiliary variables.

The implementation places the columns of $\bar{\QQ}$ in weakly
non-increasing lexicographic order. This does not affect the existence
of a factorization: columns of $\bar{\QQ}$ can be sorted by simultaneously
permuting the corresponding rows of $\HH$. Weak ordering allows equal
columns. Let $\QQ^\downarrow$ denote $\QQ$ with its columns sorted in the
same order. Among sorted candidates, equivalence up to column permutation
is exactly equality to $\QQ^\downarrow$. Non-equivalence therefore requires
only the single clause
\begin{equation}\label{eq:supp-sat-exclude-x}
\left(\bigvee_{(j,\ell):q^\downarrow_{j\ell}=0}x_{j\ell}\right)
\vee
\left(\bigvee_{(j,\ell):q^\downarrow_{j\ell}=1}\neg x_{j\ell}\right).
\end{equation}
This clause uses no auxiliary variables and has $JK$ literal occurrences.
The code factorizes $\QQ$ in its supplied column order and uses the sorted
copy only in~\eqref{eq:supp-sat-exclude-x}. This is equivalent to sorting
$\QQ$ before the construction, as in the main text, while retaining the
original column coordinates in the returned factorization.

The two-column condition ensures that every column of $\QQ$ is nonzero.
Consequently, every column of $\HH$ in a valid factorization is nonzero.
We retain these implied constraints explicitly:
\begin{equation}\label{eq:supp-sat-h-nonempty}
\bigvee_{\ell=1}^K h_{\ell k},\qquad k\in[K].
\end{equation}
The implementation also requires every row of $\bar{\QQ}$ to be nonzero,
which is implied by the factorization when the rows of $\QQ$ are nonzero.
Any optional restriction on the comparison class, such as nonzero columns
of $\bar{\QQ}$, is imposed separately.

Every nonequivalent factorization extends to a satisfying assignment of
the CNF formula after sorting the candidate columns and permuting the
corresponding rows of $\HH$: select a valid factor term for each positive
entry of $\QQ$ and assign the lexicographic auxiliary variables accordingly.
Conversely, every satisfying assignment supplies a factorization and a
sorted candidate different from $\QQ^\downarrow$, and therefore not
equivalent to $\QQ$. Under Proposition~\ref{prop:sat-nonequivalence},
satisfiability demonstrates failure of the algebraic condition, whereas
unsatisfiability establishes the condition.

The exclusion clause and the explicit nonempty-column constraints on
$\HH$ together use no auxiliary variables, $K+1$ clauses, and $JK+K^2$
literal occurrences. The Boolean-factorization clauses use
$K\|\QQ\|_0$ auxiliary variables, $JK^2+(K+1)\|\QQ\|_0$ clauses, and
$2JK^2+3K\|\QQ\|_0$ literal occurrences, where $\|\QQ\|_0$ is the
number of positive entries of $\QQ$. In particular, each positive entry
uses $2K+1$ clauses. Together with the $JK+K^2$ matrix variables and the
$O(JK)$ lexicographic encoding, the complete formulation has $O(JK^2)$
variables, clauses, and literal occurrences. These bounds describe the
encoding size, not the time required to solve the SAT instance.

The finalized default is \texttt{exclude\_x}. The earlier cardinality
encodings remain available through explicit options for reproducible
comparisons. The optional maximal-candidate restriction is off by default.

An exhaustive-search comparison can be obtained by enumerating only
$\HH$. For fixed $\HH$, define $\bar q^{\max}_{j\ell}=1$ precisely when
$h_{\ell k}=0$ for every $k$ such that $q_{jk}=0$.
Every feasible $\bar{\QQ}$ is entrywise bounded by
$\bar{\QQ}^{\max}$, and
\[
\QQ=\bar{\QQ}\odot\HH
\quad\Longrightarrow\quad
\QQ=\bar{\QQ}^{\max}\odot\HH.
\]
Indeed, increasing $\bar{\QQ}$ to $\bar{\QQ}^{\max}$ retains all
required positive entries in the product, and its definition prevents
any positive term at a zero entry of $\QQ$. This enlargement preserves
nonzero rows and nonzero columns when these are required; under
Proposition~\ref{prop:sat-nonequivalence}, the unchanged cardinality
of $\HH$ also preserves non-equivalence. Thus, for the comparison
class consisting of binary matrices with these nonzero requirements,
each of the $2^{K^2}$ assignments to $\HH$ can be checked in
$O(JK^2)$ time by constructing $\bar{\QQ}^{\max}$ and checking the
constraints. This gives the bound $O(2^{K^2}JK^2)$ for an elementary
exhaustive search. It is not a running-time prediction for the SAT
solver, and additional restrictions on the comparison class would
require a separate justification of this enlargement.

We implement the formulation using PySAT \citep{ignatiev2018pysat},
with \texttt{Glucose42} as the default solver. A satisfiable result
provides values for the entries of $\bar{\QQ}$ and $\HH$, and the
implementation independently verifies the returned factorization
and non-equivalence. An unsatisfiable result establishes that no such
assignment exists.
